% Доля давления излучения (1-beta) в центре звезды как функция lg(mu^2 M):
% решение уравнения beta^4/(1-beta) = (30.5/b0) / (mu^2 M)^2
% для однородной звезды (n=0, b0=1) и стандартной модели Эддингтона (n=3, b0=0.092).
% Кривые посчитаны точно, до 1-beta = 0.95. Компиляция: pdflatex beta-mass-figure.tex
\documentclass[tikz,border=8pt]{standalone}
\usepackage[T2A]{fontenc}
\usepackage[utf8]{inputenc}
\usepackage[russian]{babel}
\usetikzlibrary{arrows.meta}

\begin{document}
\begin{tikzpicture}[
    x=2.2cm, y=5.5cm,
    >={Stealth[length=2.4mm]},
    axis/.style = {->, semithick},
  ]

  % ================= оси =================
  \draw[axis] (-0.6,0) -- (4.45,0);
  \node at (4.08,0) [above=4pt] {$\lg(\mu^2 M)$};
  \draw[axis] (0,0) -- (0,1.08);
  \node at (0,1.06) [right=5pt] {$1-\beta$};

  \foreach \x/\t in {0/0.0, 1/1.0, 2/2.0, 3/3.0, 4/4.0} {
    \draw (\x,0) -- ++(0,-3pt) node[below=2pt] {\small \t};
  }
  \foreach \y/\t in {0.25/0.25, 0.5/0.50, 0.75/0.75, 1.0/1.00} {
    \draw (0,\y) -- ++(-3pt,0) node[left=2pt] {\small \t};
  }

  % ================= n = 0: однородная звезда, оценка сверху =================
  \draw[gray!40!black, thick] plot coordinates {
        (-0.404,0.005) (-0.299,0.008) (-0.208,0.012) (-0.128,0.017) (-0.057,0.023)
        (0.007,0.030) (0.079,0.040) (0.161,0.055) (0.228,0.070) (0.284,0.085)
        (0.334,0.100) (0.379,0.115) (0.420,0.130) (0.459,0.145) (0.496,0.160)
        (0.531,0.175) (0.565,0.190) (0.597,0.205) (0.629,0.220) (0.660,0.235)
        (0.691,0.250) (0.721,0.265) (0.751,0.280) (0.781,0.295) (0.810,0.310)
        (0.839,0.325) (0.869,0.340) (0.898,0.355) (0.928,0.370) (0.957,0.385)
        (0.987,0.400) (1.017,0.415) (1.047,0.430) (1.078,0.445) (1.109,0.460)
        (1.140,0.475) (1.172,0.490) (1.205,0.505) (1.238,0.520) (1.271,0.535)
        (1.306,0.550) (1.341,0.565) (1.377,0.580) (1.414,0.595) (1.453,0.610)
        (1.492,0.625) (1.533,0.640) (1.575,0.655) (1.618,0.670) (1.663,0.685)
        (1.710,0.700) (1.760,0.715) (1.811,0.730) (1.865,0.745) (1.922,0.760)
        (1.982,0.775) (2.047,0.790) (2.115,0.805) (2.189,0.820) (2.268,0.835)
        (2.355,0.850) (2.450,0.865) (2.556,0.880) (2.676,0.895) (2.813,0.910)
        (2.975,0.925) (3.172,0.940) (3.333,0.950)
  };
  \node[gray!40!black] at (1.02,0.60) {$n=0$};

  % ================= n = 3: стандартная модель Эддингтона =================
  \draw[blue!55!black, very thick] plot coordinates {
        (0.114,0.005) (0.219,0.008) (0.310,0.012) (0.390,0.017) (0.461,0.023)
        (0.525,0.030) (0.597,0.040) (0.680,0.055) (0.746,0.070) (0.802,0.085)
        (0.852,0.100) (0.897,0.115) (0.938,0.130) (0.977,0.145) (1.014,0.160)
        (1.049,0.175) (1.083,0.190) (1.115,0.205) (1.147,0.220) (1.178,0.235)
        (1.209,0.250) (1.239,0.265) (1.269,0.280) (1.299,0.295) (1.328,0.310)
        (1.358,0.325) (1.387,0.340) (1.416,0.355) (1.446,0.370) (1.475,0.385)
        (1.505,0.400) (1.535,0.415) (1.565,0.430) (1.596,0.445) (1.627,0.460)
        (1.658,0.475) (1.690,0.490) (1.723,0.505) (1.756,0.520) (1.790,0.535)
        (1.824,0.550) (1.859,0.565) (1.895,0.580) (1.933,0.595) (1.971,0.610)
        (2.010,0.625) (2.051,0.640) (2.093,0.655) (2.136,0.670) (2.181,0.685)
        (2.229,0.700) (2.278,0.715) (2.329,0.730) (2.383,0.745) (2.440,0.760)
        (2.501,0.775) (2.565,0.790) (2.633,0.805) (2.707,0.820) (2.786,0.835)
        (2.873,0.850) (2.968,0.865) (3.074,0.880) (3.194,0.895) (3.331,0.910)
        (3.493,0.925) (3.691,0.940) (3.851,0.950)
  };
  \node[blue!55!black] at (2.35,0.50) {$n=3$};

\end{tikzpicture}
\end{document}
